\chapter{Ejecucion, seguridad y cumplimiento normativo}

\section{Secuencia de ejecucion}
El orden referencial es: obra civil y canalizaciones enterradas, pozos de
puesta a tierra, alimentadores y subalimentadores, tableros, circuitos
derivados, luminarias y tomacorrientes, equipos de playa (surtidores, bombas
sumergibles, compresor, bombas de servicio), grupo electrogeno con ATS, UPS,
pararrayo, pruebas y verificaciones, y entrega de documentacion.

\section{Areas peligrosas}
El grifo corresponde a un lugar de manipulacion de combustibles de la Regla
120-000 del CNE-U, complementada por la Seccion 110 para lugares peligrosos.
La propuesta academica de este anteproyecto delimita zonas de riesgo alrededor
de tanques, venteos y surtidores (zonas 1 y 2) en la lamina IE-06. Los limites
definitivos y la seleccion de equipos de area clasificada deben ser revisados
por un ingeniero electricista o mecanico-electricista colegiado antes de la
construccion.

\begin{center}
\small
\begin{tabularx}{\linewidth}{@{}l c c X@{}}
\toprule
Elemento de riesgo & Zona & Radio (m) & Coordenada local (x, y) \\
\midrule
Tanque TK-1 (Gasohol Premium) & 1 & 2,2 & (14,57; 7,95) \\
Tanque TK-2 (Gasohol Regular) & 1 & 2,2 & (17,49; 8,21) \\
Tanque TK-3 (Diesel B5-S50) & 1 & 2,2 & (21,04; 8,42) \\
Tuberias de venteo & 1 & 1,5 & (0,83; 4,37) \\
Surtidor isla 1 & 2 & 3,5 & (14,44; 10,92) \\
Surtidor isla 2 & 2 & 3,5 & (19,85; 10,93) \\
\bottomrule
\end{tabularx}
\end{center}

Las coordenadas son locales (en metros, origen en la esquina suroeste del lote
a ejecutar) y corresponden a las posiciones observadas en el DWG original
(PLANO DE DISTRIBUCION.dwg). La Zona 1 se extiende 2,2~m alrededor de las bocas
de los tanques y 1,5~m alrededor de los venteos; la Zona 2 se extiende 3,5~m
alrededor de los surtidores. Estos radios son una propuesta academica sujeta a
revision competente.

\section{Normas base}
\begin{itemize}
\item Codigo Nacional de Electricidad--Utilizacion, aprobado por R.M. N.\({}^\circ\)037-2006-MEM/DM y modificado por R.M. N.\({}^\circ\)175-2008-MEM/DM.
\item Norma EM.010 del RNE, aprobada por R.M. N.\({}^\circ\)083-2019-VIVIENDA.
\item Reglamento Nacional de Edificaciones.
\item Reglamento de Seguridad para Establecimientos de Venta al Publico de Combustibles y normativa sectorial vigente de la OSINERGMIN, en lo aplicable.
\end{itemize}

\section{Reglas principales citadas}
\begin{itemize}
\item Regla 030-002: seccion minima 2,5~mm\textsuperscript{2} cobre para fuerza y alumbrado.
\item Regla 030-004 y Tabla 2: capacidad de corriente de conductores.
\item Regla 050-102: caida de tension maxima 2,5~\% por alimentador o circuito derivado y 4~\% total.
\item Regla 050-104: coordinacion \(I_b \leq I_n \leq I_z\).
\item Regla 050-210 y Tabla 14: demanda para usos no residenciales.
\item Regla 060-712: limite de puesta a tierra.
\item Tabla 16: seccion minima de enlace equipotencial.
\item EM.010 articulo 7: evaluacion de demanda.
\item EM.010 articulo 11.1: alumbrado de emergencia y senales de evacuacion.
\end{itemize}

\section{Pendientes obligatorios antes de obra}
\begin{itemize}
\item Obtener factibilidad, punto de entrega e Icc de Electro Puno (la empresa
      distribuidora se asume; el Icc de 10~kA es un supuesto de anteproyecto).
\item Sustituir familias de catalogo por placas reales.
\item Verificar cotas, alturas, PAT y areas clasificadas en campo.
\item Revision humana competente de calculos, CAD y expediente.
\end{itemize}
